The barycenter $v$ of $\sigma^n$ is the point of $\sigma^n$ all of whose barycentric coordinates are equal $(= 1/(n+1))$. The (first) barycentric subdivision $bK$ of a complex $K$ is defined inductively, as follows. (1) $bK^0 = K^0$. (2) Given $bK^i$, $bK^{i+1}$ is the union of $bK^i$ and the set of all joins $v\sigma^i$, where $v$ is the barycenter of a simplex $\sigma^{i+1}$ of $K$ and $\sigma^i \in bK^i$, $\sigma^i \subset \sigma^{i+1}$. In a Euclidean complex, the dimensions of the simplexes always form a bounded set, and so this process terminates, giving $bK$.
